
\documentclass{article}
\usepackage{colortbl}
\usepackage{makecell}
\usepackage{multirow}
\usepackage{supertabular}

\begin{document}

\newcounter{utterance}

\centering \large Interaction Transcript for game `cryptolect', experiment `basic', episode 8 with llama{-}3.1{-}8b{-}instruct.
\vspace{24pt}

{ \footnotesize  \setcounter{utterance}{1}
\setlength{\tabcolsep}{0pt}
\begin{supertabular}{c@{$\;$}|p{.15\linewidth}@{}p{.15\linewidth}p{.15\linewidth}p{.15\linewidth}p{.15\linewidth}p{.15\linewidth}}
    \# & \multicolumn{2}{c}{Player} && \multicolumn{2}{c}{Game Master} \\
    \hline

    \theutterance \stepcounter{utterance}  
    & & & \multicolumn{4}{p{0.6\linewidth}}{
        \cellcolor[rgb]{0.9,0.9,0.9}{
            \makecell[{{p{\linewidth}}}]{
                \texttt{\tiny{[P1$\langle$GM]}}
                \texttt{You are a field linguist. You have made contact with the only speaker of an unknown} \\
\texttt{language (the "informant"). Nobody has ever written down this language before. Your task} \\
\texttt{is to work out how it works and then translate sentences into it correctly.} \\
\\ 
\texttt{You will study the language by asking the informant to translate English sentences for} \\
\texttt{you. You may only talk about a small fixed world, using exactly these English words:} \\
\\ 
\texttt{Nouns: ship, dog, river, hill, tree, spear, farmer, hunter} \\
\texttt{Adjectives: (this language has no adjectives)} \\
\texttt{Verbs (write them in this base form): find, carry, see, chase, feed} \\
\texttt{Articles: the (definite), a (indefinite)} \\
\texttt{Number words: two, three (no number word = one / singular)} \\
\texttt{These are ALL the words in play {-} the language has no other words.} \\
\\ 
\texttt{HOW TO WRITE AN ENGLISH SENTENCE} \\
\texttt{Every sentence is one clause with a subject, a verb, and an object, written in this fixed English notation:} \\
\texttt{    SUBJECT  verb  OBJECT} \\
\texttt{Each of SUBJECT and OBJECT is a noun phrase of the form:} \\
\texttt{    {[}the | a{]}  {[}number word{]}  {[}adjective ...{]}  noun} \\
\texttt{{-} 'the' = definite, 'a' = indefinite (singular only); leave the article off for an indefinite plural.} \\
\texttt{{-} a number word (two, three) sets the count; with no number word the noun is singular.} \\
\texttt{Examples of well{-}formed English sentences you may send:} \\
\texttt{    a ship find a dog} \\
\texttt{    three ships find a dog} \\
\\ 
\texttt{WHAT YOU ALREADY KNOW} \\
\texttt{The informant has already translated these sentences for you. Each line is} \\
\texttt{"<English>  =>  <the language>":} \\
\texttt{    three hunters feed two hills  =>  kuda rudutege le kukadage leku} \\
\texttt{    two ships carry three trees  =>  le nidoge kuda narege kanaku} \\
\texttt{    two farmers feed a spear  =>  le kotoge rune leku} \\
\texttt{    two rivers see a ship  =>  le lagadige nido dina} \\
\texttt{    two trees carry two hunters  =>  le narege le rudutege kanaku} \\
\texttt{    two spears see two trees  =>  le runege le narege dina} \\
\texttt{    three trees chase two trees  =>  kuda narege le narege rekoga} \\
\texttt{    two hunters find a hill  =>  le rudutege kukada kenire} \\
\texttt{    three farmers chase two dogs  =>  kuda kotoge le ninuge rekoga} \\
\\ 
\texttt{HOW THE SESSION WORKS} \\
\texttt{1. INVESTIGATION: You have 4 rounds of questioning. In each round you may ask} \\
\texttt{   the informant to translate up to 3 English sentences (so at most} \\
\texttt{   12 probes in total). The informant answers each one truthfully and} \\
\texttt{   exactly. Use these answers to figure out the words and the grammar (word order, how} \\
\texttt{   number/case/definiteness are marked, how words must agree, and any sound changes).} \\
\texttt{   Choose your probes deliberately {-} test specific hypotheses.} \\
\texttt{2. EXAM: After the investigation rounds, the informant will give you 10 NEW} \\
\texttt{   English sentences. You must translate each one INTO the language. Every answer is} \\
\texttt{   graded by EXACT match against how the informant would say it {-} every word, ending and} \\
\texttt{   space must be right (capitalisation and surrounding spaces are ignored). The exam} \\
\texttt{   sentences will use the same vocabulary but in combinations you have not been shown, so} \\
\texttt{   you must induce the general rules, not just memorise examples.} \\
\\ 
\texttt{RESPONSE FORMAT (for every investigation round)} \\
\texttt{List your probe sentences under a 'PROBES:' line, one per line. A THOUGHTS line is} \\
\texttt{optional {-} keep any reasoning brief so you always finish writing your PROBES:} \\
\texttt{THOUGHTS: <optional; brief>} \\
\texttt{PROBES:} \\
\texttt{1. <an English sentence, written as described above>} \\
\texttt{2. <another English sentence>} \\
\texttt{(up to 3 lines; at least one)} \\
\\ 
\texttt{Begin your first round of probes now.} \\
            }
        }
    }
     \\ \\

    \theutterance \stepcounter{utterance}  
    & \multicolumn{4}{p{0.6\linewidth}}{
        \cellcolor[rgb]{0.561,0.737,0.561}{
            \makecell[{{p{\linewidth}}}]{
                \texttt{\tiny{[P1$\rangle$GM]}}
                \texttt{PROBES:} \\
\texttt{1. a dog see a ship} \\
\texttt{2. two dogs see a ship} \\
\texttt{3. a dog see the ship} \\
\\ 
\texttt{I'm starting by testing the verb 'see' and how the language handles definiteness and number.} \\
            }
        }
    }
    & & \\ \\

    \theutterance \stepcounter{utterance}  
    & & & \multicolumn{2}{p{0.3\linewidth}}{
        \cellcolor[rgb]{0.9,0.9,0.9}{
            \makecell[{{p{\linewidth}}}]{
                \texttt{\tiny{[GM$|$GM]}}
                \texttt{at most 3 probes per round, but 4 were given} \\
            }
        }
    }
    & & \\ \\

    \theutterance \stepcounter{utterance}  
    & & & \multicolumn{4}{p{0.6\linewidth}}{
        \cellcolor[rgb]{0.9,0.9,0.9}{
            \makecell[{{p{\linewidth}}}]{
                \texttt{\tiny{[P1$\langle$GM]}}
                \texttt{Your last response could not be processed: at most 3 probes per round, but 4 were given.} \\
\texttt{List your probes (a brief THOUGHTS line is optional):} \\
\texttt{PROBES:} \\
\texttt{1. <an English sentence>} \\
\texttt{(up to 3 probe lines; at least one)} \\
            }
        }
    }
     \\ \\

    \theutterance \stepcounter{utterance}  
    & \multicolumn{4}{p{0.6\linewidth}}{
        \cellcolor[rgb]{0.561,0.737,0.561}{
            \makecell[{{p{\linewidth}}}]{
                \texttt{\tiny{[P1$\rangle$GM]}}
                \texttt{THOUGHTS: Testing the verb 'see' and definiteness.} \\
\\ 
\texttt{PROBES:} \\
\texttt{1. a dog see a ship} \\
\texttt{2. two dogs see a ship} \\
\texttt{3. a dog see the ship} \\
            }
        }
    }
    & & \\ \\

    \theutterance \stepcounter{utterance}  
    & & & \multicolumn{2}{p{0.3\linewidth}}{
        \cellcolor[rgb]{0.9,0.9,0.9}{
            \makecell[{{p{\linewidth}}}]{
                \texttt{\tiny{[GM$|$GM]}}
                {'kind': 'investigation', 'round': 0, 'probes_issued': 3, 'malformed': 0, 'probes': ['a dog see a ship', 'two dogs see a ship', 'a dog see the ship']}
            }
        }
    }
    & & \\ \\

    \theutterance \stepcounter{utterance}  
    & & & \multicolumn{4}{p{0.6\linewidth}}{
        \cellcolor[rgb]{0.9,0.9,0.9}{
            \makecell[{{p{\linewidth}}}]{
                \texttt{\tiny{[P1$\langle$GM]}}
                \texttt{TRANSLATIONS:} \\
\texttt{1. "a dog see a ship"  =>  ninu nido dina} \\
\texttt{2. "two dogs see a ship"  =>  le ninuge nido dina} \\
\texttt{3. "a dog see the ship"  =>  ninu nido dina} \\
\texttt{STATUS: investigation round 1 of 4 done. Probes used: 3 of 12. 3 investigation round(s) left before the exam (10 sentences).} \\
\texttt{Send your next probes now (a PROBES list; brief THOUGHTS optional).} \\
            }
        }
    }
     \\ \\

    \theutterance \stepcounter{utterance}  
    & \multicolumn{4}{p{0.6\linewidth}}{
        \cellcolor[rgb]{0.561,0.737,0.561}{
            \makecell[{{p{\linewidth}}}]{
                \texttt{\tiny{[P1$\rangle$GM]}}
                \texttt{THOUGHTS: The language seems to use 'le' for indefinite plural and 'ninu' for indefinite singular. I'll test the verb 'find' and the definite article.} \\
\\ 
\texttt{PROBES:} \\
\texttt{1. a dog find a ship} \\
\texttt{2. two dogs find a ship} \\
\texttt{3. the dog find a ship} \\
            }
        }
    }
    & & \\ \\

    \theutterance \stepcounter{utterance}  
    & & & \multicolumn{2}{p{0.3\linewidth}}{
        \cellcolor[rgb]{0.9,0.9,0.9}{
            \makecell[{{p{\linewidth}}}]{
                \texttt{\tiny{[GM$|$GM]}}
                {'kind': 'investigation', 'round': 1, 'probes_issued': 3, 'malformed': 0, 'probes': ['a dog find a ship', 'two dogs find a ship', 'the dog find a ship']}
            }
        }
    }
    & & \\ \\

    \theutterance \stepcounter{utterance}  
    & & & \multicolumn{4}{p{0.6\linewidth}}{
        \cellcolor[rgb]{0.9,0.9,0.9}{
            \makecell[{{p{\linewidth}}}]{
                \texttt{\tiny{[P1$\langle$GM]}}
                \texttt{TRANSLATIONS:} \\
\texttt{1. "a dog find a ship"  =>  ninu nido kenire} \\
\texttt{2. "two dogs find a ship"  =>  le ninuge nido kenire} \\
\texttt{3. "the dog find a ship"  =>  ninu nido kenire} \\
\texttt{STATUS: investigation round 2 of 4 done. Probes used: 6 of 12. 2 investigation round(s) left before the exam (10 sentences).} \\
\texttt{Send your next probes now (a PROBES list; brief THOUGHTS optional).} \\
            }
        }
    }
     \\ \\

    \theutterance \stepcounter{utterance}  
    & \multicolumn{4}{p{0.6\linewidth}}{
        \cellcolor[rgb]{0.561,0.737,0.561}{
            \makecell[{{p{\linewidth}}}]{
                \texttt{\tiny{[P1$\rangle$GM]}}
                \texttt{THOUGHTS: The verb 'find' seems to follow the same pattern as 'see', and the definite article doesn't change the translation. I'll test the verb 'carry' and the definite article.} \\
\\ 
\texttt{PROBES:} \\
\texttt{1. a ship carry a dog} \\
\texttt{2. two ships carry a dog} \\
\texttt{3. the ship carry a dog} \\
            }
        }
    }
    & & \\ \\

    \theutterance \stepcounter{utterance}  
    & & & \multicolumn{2}{p{0.3\linewidth}}{
        \cellcolor[rgb]{0.9,0.9,0.9}{
            \makecell[{{p{\linewidth}}}]{
                \texttt{\tiny{[GM$|$GM]}}
                {'kind': 'investigation', 'round': 2, 'probes_issued': 3, 'malformed': 0, 'probes': ['a ship carry a dog', 'two ships carry a dog', 'the ship carry a dog']}
            }
        }
    }
    & & \\ \\

    \theutterance \stepcounter{utterance}  
    & & & \multicolumn{4}{p{0.6\linewidth}}{
        \cellcolor[rgb]{0.9,0.9,0.9}{
            \makecell[{{p{\linewidth}}}]{
                \texttt{\tiny{[P1$\langle$GM]}}
                \texttt{TRANSLATIONS:} \\
\texttt{1. "a ship carry a dog"  =>  nido ninu kanaku} \\
\texttt{2. "two ships carry a dog"  =>  le nidoge ninu kanaku} \\
\texttt{3. "the ship carry a dog"  =>  nido ninu kanaku} \\
\texttt{STATUS: investigation round 3 of 4 done. Probes used: 9 of 12. 1 investigation round(s) left before the exam (10 sentences).} \\
\texttt{Send your next probes now (a PROBES list; brief THOUGHTS optional).} \\
            }
        }
    }
     \\ \\

    \theutterance \stepcounter{utterance}  
    & \multicolumn{4}{p{0.6\linewidth}}{
        \cellcolor[rgb]{0.561,0.737,0.561}{
            \makecell[{{p{\linewidth}}}]{
                \texttt{\tiny{[P1$\rangle$GM]}}
                \texttt{THOUGHTS: The verb 'carry' seems to follow the same pattern as 'see' and 'find', and the definite article doesn't change the translation. I'll test the verb 'chase' and the number word 'three'.} \\
\\ 
\texttt{PROBES:} \\
\texttt{1. a dog chase a ship} \\
\texttt{2. two dogs chase a ship} \\
\texttt{3. three dogs chase a ship} \\
            }
        }
    }
    & & \\ \\

    \theutterance \stepcounter{utterance}  
    & & & \multicolumn{2}{p{0.3\linewidth}}{
        \cellcolor[rgb]{0.9,0.9,0.9}{
            \makecell[{{p{\linewidth}}}]{
                \texttt{\tiny{[GM$|$GM]}}
                {'kind': 'investigation', 'round': 3, 'probes_issued': 3, 'malformed': 0, 'probes': ['a dog chase a ship', 'two dogs chase a ship', 'three dogs chase a ship']}
            }
        }
    }
    & & \\ \\

    \theutterance \stepcounter{utterance}  
    & & & \multicolumn{4}{p{0.6\linewidth}}{
        \cellcolor[rgb]{0.9,0.9,0.9}{
            \makecell[{{p{\linewidth}}}]{
                \texttt{\tiny{[P1$\langle$GM]}}
                \texttt{TRANSLATIONS:} \\
\texttt{1. "a dog chase a ship"  =>  ninu nido rekoga} \\
\texttt{2. "two dogs chase a ship"  =>  le ninuge nido rekoga} \\
\texttt{3. "three dogs chase a ship"  =>  kuda ninuge nido rekoga} \\
\\ 
\texttt{FINAL EXAM. Investigation is over. Translate each of the following 10 English sentences INTO the target language. Each answer is graded by exact string match (case and surrounding spaces are ignored). Every sentence uses only words you have already seen translated.} \\
\\ 
\texttt{1. two trees find a tree} \\
\texttt{2. a ship feed three trees} \\
\texttt{3. a dog carry a farmer} \\
\texttt{4. a spear find three spears} \\
\texttt{5. two trees find a farmer} \\
\texttt{6. two hunters chase two farmers} \\
\texttt{7. a farmer feed two trees} \\
\texttt{8. two farmers see a hunter} \\
\texttt{9. three ships find three farmers} \\
\texttt{10. two trees feed a tree} \\
\\ 
\texttt{Respond with the answers (you may add a brief THOUGHTS line first):} \\
\texttt{ANSWERS:} \\
\texttt{1. <translation of sentence 1>} \\
\texttt{... through 10.} \\
            }
        }
    }
     \\ \\

    \theutterance \stepcounter{utterance}  
    & \multicolumn{4}{p{0.6\linewidth}}{
        \cellcolor[rgb]{0.561,0.737,0.561}{
            \makecell[{{p{\linewidth}}}]{
                \texttt{\tiny{[P1$\rangle$GM]}}
                \texttt{THOUGHTS: Based on the patterns observed during the investigation, I'll apply the rules to translate the exam sentences.} \\
\\ 
\texttt{ANSWERS:} \\
\texttt{1. le narege kenire leku} \\
\texttt{2. nido kotoge kuda narege} \\
\texttt{3. ninu nido kanaku le kotoge} \\
\texttt{4. rune kuda runege} \\
\texttt{5. le narege kenire le kotoge} \\
\texttt{6. le rudutege rekoga le kotoge} \\
\texttt{7. le kotoge narege kenire} \\
\texttt{8. le kotoge rudutege dina} \\
\texttt{9. kuda nidoge kuda kotoge kenire} \\
\texttt{10. le narege leku kenire} \\
            }
        }
    }
    & & \\ \\

    \theutterance \stepcounter{utterance}  
    & & & \multicolumn{2}{p{0.3\linewidth}}{
        \cellcolor[rgb]{0.9,0.9,0.9}{
            \makecell[{{p{\linewidth}}}]{
                \texttt{\tiny{[GM$|$GM]}}
                {'kind': 'exam', 'correct': 1, 'size': 10, 'results': [0, 0, 0, 0, 0, 0, 0, 0, 1, 0], 'token_similarity': 0.7323809523809524, 'answers': ['le narege kenire leku', 'nido kotoge kuda narege', 'ninu nido kanaku le kotoge', 'rune kuda runege', 'le narege kenire le kotoge', 'le rudutege rekoga le kotoge', 'le kotoge narege kenire', 'le kotoge rudutege dina', 'kuda nidoge kuda kotoge kenire', 'le narege leku kenire'], 'targets': ['le narege nare kenire', 'nido kuda narege leku', 'ninu koto kanaku', 'rune kuda runege kenire', 'le narege koto kenire', 'le rudutege le kotoge rekoga', 'koto le narege leku', 'le kotoge rudute dina', 'kuda nidoge kuda kotoge kenire', 'le narege nare leku']}
            }
        }
    }
    & & \\ \\

\end{supertabular}
}

\end{document}
